% fig7_emulator_levers_and_horizon.tex -- where the emulator bottleneck is NOT, and why speed
% is irrelevant.
%
% COMPILE
%   pdflatex -interaction=nonstopmode fig7_emulator_levers_and_horizon.tex
%   dvisvgm --pdf --no-fonts --exact fig7_emulator_levers_and_horizon.pdf
%
% WHY PANEL a IS A TABLE AND NOT A BAR CHART
% The six levers are measured against DIFFERENT baselines, on different cadences and different
% tracer sets: +3.75 is a gap closure in log space on daily eqpac, +0.14 is a linear-skill delta
% on the same domain, +0.007 is a capacity sweep on monthly 3-D. Putting them on one axis would
% imply they are commensurable and would make log space look ~500x the lever that ensembling is,
% which is an artifact of the units rather than a finding. A table with a verdict stripe states
% each effect in its own units and lets the VERDICT carry the comparison.
%
% SEMANTICS OF THE COLOUR LADDER HERE (restated, as in fig6 -- this is not a parameter figure)
%   ddBlue  load-bearing: removing it measurably degrades the result
%   ddGold  not a lever: measured, and the effect is at or near zero
%   ddRed   STROKE ONLY -- the retracted curve in panel b, i.e. a number that was withdrawn
%
% SOURCES
%   docs/findings/2026-07-19_emulator_honest_bounds.md   both panels (lever table; k8 rollout)
%   docs/findings/2026-07-30_daily_logspace_training.md  the +3.75 gap closure in panel a
\documentclass[border=5pt]{standalone}
\usepackage{darwindiff-tikz-v2}

\begin{document}
\begin{tikzpicture}[
  ddval/.style  = {ddnum,  anchor=west},
  ddcap/.style  = {ddnote, anchor=west, text width=76mm, align=left},
]

%% =====================================================================================
%% PANEL a -- the lever ladder
%% =====================================================================================
\begin{scope}[shift={(0,0)}]
  \node[ddpanel, anchor=west] at (-0.15, 1.30) {a};
  \node[ddbody, anchor=west, font=\sffamily\fHead] at (0.30, 1.30)
       {Three levers are real; three are measured and dead};

  \node[ddhead, anchor=west] at (0.30, 0.86) {Lever};
  \node[ddhead, anchor=west] at (4.05, 0.86) {Measured effect};
  \draw[ddhair] (0.00, 0.64) -- (8.20, 0.64);

  \node[ddhead, anchor=west, text=ddBlue] at (0.30, 0.34) {LOAD-BEARING};
  \foreach \i/\nm/\eff in {
    0/{log-space transform}/{closes 91.6\% of the gap; 0\% negatives by construction},
    1/{deep ensembling, 8 seeds}/{$+0.14$ daily, $+0.05$ 3-D; 3$\times$ fewer negatives},
    2/{rollout-aware training (k8)}/{STABILITY ONLY: mass 1.000 vs k1 diverging to 3.05e8}}{
      \pgfmathsetmacro{\yy}{-0.02 - \i*0.36}
      \draw[ddBlue, line width=2.2pt] (0.00, \yy) -- (0.18, \yy);
      \node[ddbody, anchor=west] at (0.30, \yy) {\nm};
      \node[ddnote, anchor=west] at (4.05, \yy) {\eff};
  }

  \node[ddhead, anchor=west, text=ddGold] at (0.30, -1.20) {NOT A LEVER};
  \foreach \i/\nm/\eff in {
    0/{capacity, up to ${\sim}4\times$ params}/{$+0.007$ \textemdash{} saturated},
    1/{EDM diffusion corrector}/{zero skill, $-0.03$ on 3-D, at 2048 passes vs 8},
    2/{architectural diversity}/{calibration 0.240, unchanged}}{
      \pgfmathsetmacro{\yy}{-1.56 - \i*0.36}
      \draw[ddGold, line width=2.2pt] (0.00, \yy) -- (0.18, \yy);
      \node[ddbody, anchor=west] at (0.30, \yy) {\nm};
      \node[ddnote, anchor=west] at (4.05, \yy) {\eff};
  }

  \node[ddcap, anchor=north west] at (0.00, -2.85)
    {Effects are quoted in their own units because they are measured against different
     baselines; the stripe, not a shared axis, carries the comparison. The pattern is the
     finding: \emph{capacity and architecture are not the bottleneck}. Roughly $4\times$ the
     parameters bought $+0.007$, and a diffusion corrector running 2048 forward passes per
     field bought nothing a plain 8-member ensemble did not already give.};
\end{scope}

%% =====================================================================================
%% PANEL b -- the horizon is one step, and the longer one was a calendar artifact
%% =====================================================================================
\begin{scope}[shift={(9.80,0)}]
  \node[ddpanel, anchor=west] at (-0.15, 1.30) {b};
  \node[ddbody, anchor=west, font=\sffamily\fHead] at (0.30, 1.30)
       {The horizon is one step, not nine months};

  \node[ddhead, anchor=west] at (0.30, 0.86) {Skill vs a correctly-binned seasonal climatology};
  \draw[ddhair] (0.00, 0.64) -- (8.20, 0.64);

  % axes: x = rollout step 1..12 over 5.6cm from x=0.80 ; y = -0.80..+0.70 mapped into the band
  % y in [-1.30, +0.40], i.e. BELOW the header rule at 0.64. Getting this wrong once put the
  % whole plot on top of the panel title.
  \pgfmathsetmacro{\ox}{0.80}
  \pgfmathsetmacro{\sx}{5.60/11.0}
  \pgfmathsetmacro{\oy}{-1.30}
  \pgfmathsetmacro{\sy}{1.70/1.50}

  % y gridlines
  \foreach \v in {-0.75,-0.50,-0.25,0,0.25,0.50}{
    \pgfmathsetmacro{\yy}{\oy + (\v+0.80)*\sy}
    \draw[ddRule, line width=0.5pt] (\ox-0.15, \yy) -- (\ox+5.75, \yy);
    \node[ddnote, anchor=east] at (\ox-0.22, \yy) {\v};
  }
  % zero line -- the climatology itself
  \pgfmathsetmacro{\yzero}{\oy + 0.80*\sy}
  \draw[ddInk, line width=0.9pt] (\ox-0.15, \yzero) -- (\ox+5.75, \yzero);
  % NOT labelled: the panel subtitle already names the baseline, and any label on this line
  % lands on either the retracted curve or the -0.25 gridline.

  % x ticks
  \foreach \s in {1,2,3,6,9,12}{
    \pgfmathsetmacro{\xx}{\ox + (\s-1)*\sx}
    \node[ddnote, anchor=north] at (\xx, \oy-0.08) {\s};
  }
  \node[ddnote, anchor=north] at ({\ox+2.80}, \oy-0.40) {rollout step \ (1 step ${\approx}$ 2 months)};

  % OLD, buggy bins -- RETRACTED. Red, stroke only, dashed: a withdrawn number.
  \draw[ddRed, line width=0.9pt, dash pattern=on 3pt off 2pt]
    plot coordinates {
      ({\ox+0*\sx},{\oy+(0.614+0.80)*\sy})
      ({\ox+1*\sx},{\oy+(0.469+0.80)*\sy})
      ({\ox+2*\sx},{\oy+(0.432+0.80)*\sy})
      ({\ox+5*\sx},{\oy+(0.259+0.80)*\sy})
      ({\ox+8*\sx},{\oy+(0.253+0.80)*\sy})
      ({\ox+11*\sx},{\oy+(0.028+0.80)*\sy})};

  % NEW, correct bins -- the honest curve
  \draw[ddBlue, line width=1.3pt]
    plot coordinates {
      ({\ox+0*\sx},{\oy+(0.240+0.80)*\sy})
      ({\ox+1*\sx},{\oy+(-0.018+0.80)*\sy})
      ({\ox+2*\sx},{\oy+(-0.012+0.80)*\sy})
      ({\ox+5*\sx},{\oy+(-0.420+0.80)*\sy})
      ({\ox+8*\sx},{\oy+(-0.292+0.80)*\sy})
      ({\ox+11*\sx},{\oy+(-0.748+0.80)*\sy})};
  \foreach \i/\v in {0/0.240, 1/-0.018, 2/-0.012, 5/-0.420, 8/-0.292, 11/-0.748}{
    \filldraw[ddBlue] ({\ox+\i*\sx},{\oy+(\v+0.80)*\sy}) circle (1.6pt);
  }

  % the crossing
  \pgfmathsetmacro{\xcross}{\ox + 1*\sx}
  \draw[ddMute, line width=0.6pt, dash pattern=on 2pt off 1.5pt]
    (\xcross, \yzero) -- (\xcross, 0.44);
  \node[ddnote, anchor=south west] at (\xcross+0.06, 0.26) {below climatology};
  \node[ddnote, anchor=south west] at (\xcross+0.06, 0.08) {from step 2 onward};

  % legend, below the x-axis label
  \draw[ddRed, line width=0.9pt, dash pattern=on 3pt off 2pt] (0.30, -2.26) -- (0.62, -2.26);
  \node[ddnote, anchor=west] at (0.70, -2.26) {RETRACTED (buggy month bins)};
  \draw[ddBlue, line width=1.3pt] (4.30, -2.26) -- (4.62, -2.26);
  \node[ddnote, anchor=west] at (4.70, -2.26) {corrected};

  \node[ddcap, anchor=north west] at (0.00, -2.54)
    {The stored cube's time axis was 0.75$\times$ truth, so 94\% of climatology month-bins were
     wrong and the baseline was weakened at exactly the lags being claimed. Re-binned, k8 beats
     climatology at \emph{one step} and is at or below it thereafter. Mass is conserved to 1.000
     at every lag, so this is not blow-up \textemdash{} it is a correct-looking trajectory with
     no skill, which is why speed (7.45\,ms per step, a century in 9\,s) buys nothing.};
\end{scope}

\end{tikzpicture}
\end{document}
